% sec-03: Set A, applications 01 to 05.  Figures 5 and 8.
\section{Set A: The Surgical Perspective}

\draftinstr{full-apply writes 6,000 to 7,000 characters from
\appfile{applications/app-01-nih-pioneer-award/} through
\appfile{applications/app-05-nih-sbir-seed/}, one subsection per application,
each summarising the mechanism fit and the one argument unique to that
recipient. Do not repeat the trial description; it is given once in \S5.}

\subsection{What Leads in Set A}

\draftinstr{One paragraph: the operation leads, daraxonrasib enters as the
perioperative arm the operation is timed around, and the endpoint table is led
by conversion, anastomotic integrity, operative time, blood loss, R0 margin and
90-day morbidity.}

\figslot{plantuml-type / use case}{7.0}{Seven actors and the eleven actions each
may take. Every action has at least one human association; the advisory model
has exactly one, and the struck links are the four it is denied by construction.}

\subsection{Applications 01 to 05, One Paragraph Each}

\draftinstr{Five short subsections. 01 Pioneer Award: the unit of risk is the
individual and the bet does not survive decomposition into R01 aims. 02 ARPA-H:
three gates with numeric kill criteria. 03 NSF TIP X-Labs: five activities that
no existing performer type covers. 04 DOE Genesis Mission: the Robotics mission
sentence applied to an operating room. 05 NIH SEED SBIR: the product is a
site-portable governance package.}

\figslot{mermaid-type / flowchart with decisions}{5.2}{The five questions a
reviewer asks in order, and the section of the application that answers each. A
no at gate 1 or gate 2 ends the review; a no at gates 3 to 5 changes the ask.}

\draftinstr{full-apply builds Table 3 here: figure inventory for Set A, listing
which of the five diagram types each application used and why, sourced from the
Figures section of each \appfile{applications/app-0*/README.md}.}
